%=============================================================================
% BINARY FORK MATHEMATICS: Exact Decomposition of the Forced Microsector Fork
% in the DFD Derivation of alpha^{-1}
% Created 2026-03-22
%=============================================================================
\documentclass[12pt,a4paper]{article}
\usepackage{amsmath,amssymb,amsthm}
\usepackage{booktabs}
\usepackage{geometry}
\geometry{margin=2.5cm}
\usepackage{hyperref}
\usepackage{tcolorbox}

\newtheorem{theorem}{Theorem}
\newtheorem{proposition}{Proposition}
\newtheorem{lemma}{Lemma}

\title{Binary Fork Mathematics:\\[6pt]
\large Exact Decomposition of the Forced Microsector Fork\\
in the DFD $\alpha^{-1}$ Derivation}
\author{Analysis extracted from G.\ Alcock, \textit{Density Field Dynamics:
A Complete Unified Theory}, v108}
\date{March 2026}

\begin{document}
\maketitle

\begin{abstract}
We derive the exact mathematical relationship between the two branches
of the forced binary fork in the DFD derivation of $\alpha^{-1}$.
The key result is that the microsector trace choice (regular-module vs.\
fermion-rep) affects only the fraction $w = 7/80$ of the total
$U(1)$ stiffness coefficient that passes through the microsector trace.
This yields the exact decomposition
$\alpha^{-1} = \alpha^{-1}_{\rm base}\,[1 + w\,(\varepsilon_{\rm adj} - 1)]$,
where $\alpha^{-1}_{\rm base} = 137.033\,071\,79\ldots$ is the value at
unit trace normalization. The coefficient $G$ in the user's trial ansatz
$\alpha^{-1} = F \pm G/4095$ is \textbf{not} exactly~12; it equals
$\alpha^{-1}_{\rm base} \times w \approx 11.990$, reflecting the
hypercharge weighting of the microsector.
\end{abstract}

\tableofcontents
\newpage

%=============================================================================
\section{Setup and Notation}
%=============================================================================

From Section~8C of the Unified Theory (``Convention-Locked $\alpha$ from
the Microsector''), the fine-structure constant is determined by
\begin{equation}\label{eq:master}
\alpha^{-1} = 6\pi\,\kappa_{U(1)}\,,
\end{equation}
where $\kappa_{U(1)}$ is the renormalized $U(1)$ stiffness coefficient
computed from the spectral action on the Toeplitz-truncated
$\mathbb{C}P^2 \times S^3$ internal geometry.

The forced binary fork arises because one factor in the $a_4$
Seeley--DeWitt coefficient---the microsector trace normalization---admits
exactly two choices:

\begin{center}
\renewcommand{\arraystretch}{1.3}
\begin{tabular}{lcc}
\toprule
Branch & Hilbert space & $\varepsilon_{\rm adj}$ \\
\midrule
A (regular-module) & $\mathcal{H}_F = M_d(\mathbb{C})$
  & $\dfrac{d^2}{d^2-1} = \dfrac{4096}{4095}$ \\[8pt]
B (fermion-rep) & $\mathcal{H}_F = \mathbb{C}^d$
  & $\dfrac{d^2-1}{d^2} = \dfrac{4095}{4096}$ \\
\bottomrule
\end{tabular}
\end{center}

\noindent with $d = k_{\max} + 4 = 64$.

The reported numerical values are:
\begin{align}
\alpha^{-1}_A &= 137.035\,999\,85 && \text{(Branch A, residual $-0.006$~ppm)} \\
\alpha^{-1}_B &= 137.030\,144\,45 && \text{(Branch B, residual $+42.7$~ppm)} \\
\alpha^{-1}_{\rm exp} &= 137.035\,999\,084 && \text{(CODATA 2022)}
\end{align}

%=============================================================================
\section{The Exact Linear Decomposition}
\label{sec:linear}
%=============================================================================

\begin{proposition}[Linear dependence on $\varepsilon_{\rm adj}$]
The microsector trace normalization factor $\varepsilon_{\rm adj}$ enters
the $a_4$ Seeley--DeWitt coefficient linearly. Only the fraction
$w = N_{\rm species}/(g_F \cdot \mathrm{Tr}(Y^2)) = 7/80$ of the total
gauge-kinetic coefficient depends on the microsector trace.
Therefore:
\begin{equation}\label{eq:linear}
\boxed{
\alpha^{-1}(\varepsilon) = \alpha^{-1}_{\rm base}\,
\bigl[(1 - w) + w\,\varepsilon_{\rm adj}\bigr]
= \alpha^{-1}_{\rm base}\,
\bigl[1 + w\,(\varepsilon_{\rm adj} - 1)\bigr]
}
\end{equation}
where $\alpha^{-1}_{\rm base} \equiv \alpha^{-1}(\varepsilon = 1)$ is the
value with unit trace normalization.
\end{proposition}

\paragraph{Verification.}
From the two reported branch values and the two known $\varepsilon$ values,
we solve for the two unknowns $C_0$ and $C_1$ in the linear model
$\alpha^{-1} = C_0 + C_1\,\varepsilon$:
\begin{align}
C_1 &= \frac{\alpha^{-1}_A - \alpha^{-1}_B}
            {\varepsilon^{(A)} - \varepsilon^{(B)}}
     = \frac{0.005\,855\,40}{0.000\,488\,34}
     = 11.990\,395\ldots \\
C_0 &= \alpha^{-1}_A - C_1\,\varepsilon^{(A)}
     = 125.042\,677\ldots
\end{align}

\paragraph{Check that $C_1/C_0 = w/(1-w) = 7/73$.}
\begin{equation}
\frac{C_1}{C_0} = \frac{11.990\,395}{125.042\,677} = 0.095\,890\,41\ldots
\quad\text{vs.}\quad
\frac{7}{73} = 0.095\,890\,41\ldots
\end{equation}
Agreement to $< 0.2$~ppm, confirming the relationship $C_0 = C_1 \times 73/7$
and hence:
\begin{equation}
\alpha^{-1}_{\rm base} = C_0 + C_1
= C_1\,\frac{80}{7}
= 137.033\,071\,79\ldots
\end{equation}

%=============================================================================
\section{Why $G$ Is Not Exactly~12}
\label{sec:G-not-12}
%=============================================================================

The user's trial ansatz was:
\begin{equation}
\alpha^{-1} = F \pm \frac{G}{4095}\,,
\end{equation}
treating the fork as a symmetric perturbation with a single
$\delta\varepsilon = 1/(d^2-1) = 1/4095$.

\subsection{The asymmetry}

The fork is \textbf{not symmetric}:
\begin{align}
\varepsilon^{(A)} - 1 &= +\frac{1}{d^2-1} = +\frac{1}{4095} \\
\varepsilon^{(B)} - 1 &= -\frac{1}{d^2\phantom{{}^{-1}}} = -\frac{1}{4096}
\end{align}
Branch~A adds $1/4095$; Branch~B subtracts $1/4096$. The shifts differ by
$1/(d^2(d^2{-}1)) = 1/16\,773\,120$.

\subsection{The coefficient $G$}

In the user's symmetric ansatz, the ``$G$'' would be:
\begin{equation}
G_{\rm user} = \frac{\alpha^{-1}_A - \alpha^{-1}_B}{2}
\times (d^2-1) = 0.002\,927\,7 \times 4095 = 11.989\ldots
\end{equation}

This is \textbf{close to but not exactly~12}. The exact coefficient is:
\begin{equation}
\boxed{G = C_1 = \alpha^{-1}_{\rm base} \times w
= \alpha^{-1}_{\rm base} \times \frac{7}{80}
= 11.990\,39\ldots}
\end{equation}

The deviation from~12 is:
\begin{equation}
12 - G = 12 - 11.990 = 0.010\,,\quad \frac{12 - G}{G} = 0.08\%\,.
\end{equation}

\paragraph{Physical interpretation.}
The coefficient $G$ is not a simple integer. It equals
$\alpha^{-1}_{\rm base} \times w$, where:
\begin{itemize}
\item $\alpha^{-1}_{\rm base} = 137.033\ldots$ is itself a \emph{computed}
  quantity from the spectral action at unit trace normalization,
\item $w = 7/80 = N_{\rm species}/(g_F \cdot \mathrm{Tr}(Y^2))$ is the
  SM hypercharge weighting---the fraction of the gauge kinetic term that
  flows through the microsector trace.
\end{itemize}
The product $137.033 \times 7/80 = 11.990$ happens to be near~12 but has no
reason to equal it exactly.

\subsection{If $G$ were exactly~12}

Setting $C_1 = 12$ and using $\alpha^{-1}_A = 137.035\,999\,85$ to
determine $C_0$:
\begin{equation}
\alpha^{-1}_B\big|_{C_1=12}
= C_0 + 12 \times \frac{4095}{4096}
= 137.030\,139\,76\,,
\end{equation}
which differs from the reported $137.030\,144\,45$ by $+4.7 \times 10^{-6}$
($\sim 0.03$~ppm). This is small but real---the reported values have
$\sim 0.01$~ppm precision.

%=============================================================================
\section{Exact Formulas for Both Branches}
\label{sec:exact}
%=============================================================================

\begin{tcolorbox}[colback=blue!5!white, colframe=blue!50!black,
  title=Exact Branch Formulas]
With $w = 7/80$, $d = 64$:

\textbf{Branch A (regular-module, survives):}
\begin{equation}
\alpha^{-1}_A = \alpha^{-1}_{\rm base}\,
\left[\frac{73}{80} + \frac{7}{80}\cdot\frac{d^2}{d^2-1}\right]
= \alpha^{-1}_{\rm base}\,\frac{73(d^2{-}1) + 7d^2}{80(d^2{-}1)}
= \alpha^{-1}_{\rm base}\,\frac{80\,d^2 - 73}{80(d^2{-}1)}
\end{equation}

Numerically: $\alpha^{-1}_{\rm base} \times 327\,607/327\,600
= 137.035\,999\,85$.

\medskip

\textbf{Branch B (fermion-rep, falsified):}
\begin{equation}
\alpha^{-1}_B = \alpha^{-1}_{\rm base}\,
\left[\frac{73}{80} + \frac{7}{80}\cdot\frac{d^2-1}{d^2}\right]
= \alpha^{-1}_{\rm base}\,\frac{73\,d^2 + 7(d^2{-}1)}{80\,d^2}
= \alpha^{-1}_{\rm base}\,\frac{80\,d^2 - 7}{80\,d^2}
\end{equation}

Numerically: $\alpha^{-1}_{\rm base} \times 327\,673/327\,680
= 137.030\,144\,45$.
\end{tcolorbox}

%=============================================================================
\section{Fork Separation}
\label{sec:separation}
%=============================================================================

The exact separation between the two branches is:
\begin{align}
\alpha^{-1}_A - \alpha^{-1}_B
&= \alpha^{-1}_{\rm base}\,w\,
\left(\frac{d^2}{d^2-1} - \frac{d^2-1}{d^2}\right) \nonumber\\
&= \alpha^{-1}_{\rm base}\,w\,
\frac{d^4 - (d^2-1)^2}{d^2(d^2-1)} \nonumber\\
&= \alpha^{-1}_{\rm base}\,w\,
\frac{2d^2 - 1}{d^2(d^2-1)}
\end{align}

Numerically:
\begin{equation}
137.033 \times \frac{7}{80} \times \frac{8191}{4096 \times 4095}
= 0.005\,855\,4\,,
\end{equation}
matching $\alpha^{-1}_A - \alpha^{-1}_B = 0.005\,855\,4$ exactly.

In ppm:
\begin{equation}
\frac{\alpha^{-1}_A - \alpha^{-1}_B}{\alpha^{-1}_{\rm exp}}
= \frac{0.005855}{137.036} = 42.7\text{~ppm}\,.
\end{equation}
This is the full lever arm of the fork---exactly the 42.7~ppm residual
of Branch~B, since Branch~A sits at $-0.006$~ppm.

%=============================================================================
\section{The Midpoint vs.\ the Baseline}
\label{sec:midpoint}
%=============================================================================

The user computed the midpoint
$F = (\alpha^{-1}_A + \alpha^{-1}_B)/2 = 137.033\,072\,15$.
This is \textbf{not} the same as $\alpha^{-1}_{\rm base} = 137.033\,071\,79$
because the fork is asymmetric:

\begin{equation}
F - \alpha^{-1}_{\rm base}
= \frac{\alpha^{-1}_{\rm base}\,w}{2}
\left(\frac{1}{d^2-1} - \frac{1}{d^2}\right)
= \frac{\alpha^{-1}_{\rm base}\,w}{2\,d^2(d^2-1)}
= 3.6 \times 10^{-7}\,.
\end{equation}

The midpoint exceeds the baseline by $\sim 0.003$~ppm---negligible for
most purposes but detectable at the precision of the reported values.

%=============================================================================
\section{Summary of Results}
\label{sec:summary}
%=============================================================================

\begin{enumerate}
\item The forced binary fork acts through a \textbf{linear} dependence of
  $\alpha^{-1}$ on $\varepsilon_{\rm adj}$, with slope
  $C_1 = \alpha^{-1}_{\rm base} \times w = 11.990$.

\item The hypercharge weight $w = N_{\rm species}/(g_F \cdot \mathrm{Tr}(Y^2))
  = 7/80 = 0.0875$ determines what fraction of the gauge kinetic term
  passes through the microsector trace.

\item The coefficient $G$ in a symmetric decomposition is
  $G \approx 11.99$, \textbf{not} exactly~12.
  The near-integer coincidence arises because
  $\alpha^{-1}_{\rm base} \times 7/80 \approx 137.03 \times 0.0875
  \approx 11.99$.

\item The fork is \textbf{asymmetric}: Branch~A shifts by $+1/(d^2{-}1)$,
  Branch~B shifts by $-1/d^2$. The midpoint of the two branches differs
  from the unit-normalization baseline by $\sim 0.003$~ppm.

\item The exact formula is:
  \begin{equation}
  \alpha^{-1} = \alpha^{-1}_{\rm base}\,
  \left[\frac{73}{80} + \frac{7}{80}\,\varepsilon_{\rm adj}\right]
  \end{equation}
  where $\alpha^{-1}_{\rm base} = 6\pi\,\kappa_{U(1)}(\varepsilon{=}1)
  = 137.033\,072$ is a \emph{numerical} output of the spectral action.

\item There is \textbf{no single closed-form algebraic expression} for
  $\alpha^{-1}_{\rm base}$; it requires the full $a_4$ Seeley--DeWitt
  computation on the Toeplitz-truncated $\mathbb{C}P^2 \times S^3$.
  The approximate formula
  $\alpha^{-1} \approx (35\pi/3)\,\beta\,(d{-}1)/d\,\varepsilon_{\rm adj}$
  has an 85~ppm residual.
\end{enumerate}

%=============================================================================
\section*{Numerical Verification Script}
%=============================================================================

\begin{verbatim}
import math

# Reported branch values (from Section 8C, Table XXVII)
alpha_A = 137.03599985   # Branch A (BOOST = 4096/4095)
alpha_B = 137.03014445   # Branch B (DROP = 4095/4096)

d = 64
BOOST = d**2 / (d**2 - 1)   # 4096/4095
DROP  = (d**2 - 1) / d**2   # 4095/4096
w     = 7.0 / 80             # hypercharge weight

# Linear decomposition
C1 = (alpha_A - alpha_B) / (BOOST - DROP)
C0 = alpha_A - C1 * BOOST
alpha_base = C0 + C1

# Verify C1/C0 = w/(1-w) = 7/73
print(f"C1/C0 = {C1/C0:.10f}  vs  7/73 = {7/73:.10f}")

# Verify the formula
for label, eps in [("A", BOOST), ("B", DROP)]:
    val = alpha_base * ((1 - w) + w * eps)
    print(f"Branch {label}: {val:.10f}")

# Fork separation
sep = alpha_base * w * (2*d**2 - 1) / (d**2 * (d**2 - 1))
print(f"Fork separation: {sep:.10f}")
print(f"Actual: {alpha_A - alpha_B:.10f}")

# G coefficient
G = alpha_base * w
print(f"G = alpha_base * w = {G:.6f} (not 12)")
\end{verbatim}

Output:
\begin{verbatim}
C1/C0 = 0.0958904110  vs  7/73 = 0.0958904110
Branch A: 137.0359998497
Branch B: 137.0301444503
Fork separation: 0.0058553993
Actual: 0.0058554000
G = alpha_base * w = 11.990394 (not 12)
\end{verbatim}

\end{document}
